\documentclass[handout]{beamer}
\usepackage{graphicx}
\usepackage{xcolor}
\setbeamertemplate{blocks}[rounded][shadow=true]
\setbeamercolor{block body}{bg=gray!20, fg=black}
\setbeamercolor{block title}{bg=blue!20, fg=black}

\title{Clase 3: Sostenibilidad, Riesgo y Algoritmos}
\subtitle{De Horizontes Finitos a Decisiones de Largo Plazo}
\author{Daniela Pucci de Farias}
\institute{Universidad Torcuato Di Tella}
\date{Junio 2025}

\begin{document}

\frame{\titlepage}

\begin{frame}{Operaciones Sostenibles}
\begin{itemize}
    \item ¿Qué pasa cuando no hay una fecha de fin clara? ($T \rightarrow \infty$)
    \item \textbf{Políticas Estacionarias:} La mejor acción depende del estado, no del tiempo.
    \item \textbf{Ejemplo: Pesca de Salmones}
    \begin{itemize}
        \item Horizonte corto: Sobreexplotación para ganancias rápidas.
        \item Horizonte infinito: Gestión sostenible para garantizar el futuro.
    \end{itemize}
\end{itemize}
\end{frame}

\begin{frame}{Algoritmos: ¿Cómo decide una IA?}
\begin{itemize}
    \item \textbf{Iteración de Valor:} Refinar el "tablero de puntajes" hasta que converja.
    \item \textbf{Iteración de Política:} Evaluar una estrategia actual y buscar una mejora incremental.
    \item \textbf{Rollout (Mejora de Heurísticas):}
    \begin{itemize}
        \item Tomar una regla de negocio actual y usarla para simular un paso adelante.
        \item Base de aplicaciones exitosas como AlphaGo.
    \end{itemize}
\end{itemize}
\end{frame}

\begin{frame}{Riesgo y Utilidad en los Negocios}
\begin{itemize}
    \item Maximizar el valor esperado no siempre es el objetivo real.
    \item \textbf{Aversión al Riesgo:} Preferimos \$500 seguros a una apuesta de \$0 o \$1000.
    \item \textbf{Función de Utilidad $u(w)$:}
    \begin{itemize}
        \item Captura la psicología del tomador de decisiones.
        \item Ejemplo: $u(w) = \log(w)$ (Utilidad logarítmica).
    \end{itemize}
\end{itemize}
\end{frame}

\begin{frame}{Optimización de Portafolio}
\begin{itemize}
    \item Aplicación de MDP en finanzas:
    \begin{itemize}
        \item \textbf{Acción:} ¿Qué porcentaje invertir en activos riesgosos?
        \item \textbf{Estado:} Riqueza acumulada.
    \end{itemize}
    \item \textbf{Conclusión:} La política óptima depende de la volatilidad y de nuestra propia aversión al riesgo definida en el MDP.
\end{itemize}
\end{frame}

\end{document}
